%\chapter{Xa: Curved Plates or Shells}
% 235 b  395 (409/495)
235 b. In the last chapter 2: 2, 238) we
have considered the comparatively simple problem of the vibration in two
dimensions of a cylindrical shell, so far at least as relates to vibrations of
a flexural character. The shell is supposed to be thin, to be composed of
isotropic material, and to be bounded by infinite coaxial cylindrical surfaces.
It is proposed in the present chapter to treat the problem of the cylindrical
shell more generally, and further to give the theory of the flexural vibrations
of spherical shells.

In considering the deformation of a thin shell the most important question
which presents itself is whether the middle surface, viz. the surface which
lies midway between the boundaries, does, or does not, undergo extension. In
the former case the deformation may be called extensional, and its potential
energy is proportional to the thickness of the shell, which will be denoted by
\(2h\). Since the inertia of the shell, and therefore the kinetic energy of a given
motion, is also proportional to , the frequencies of vibration are in this
case independent of r, § 44. On the other hand, when no line traced upon the
middle surface undergoes extension, the potential energy of a deformation is of
a higher order in the small quantity h. If the shell be conceived to be divided
into laminae, the extension in any lamina is proportional to its distance from
the middle surface, and the contribution to the potential energy is
proportional to the square of that distance. When the integration over the
thickness is carried out, the whole potential energy is found to be
proportional to h Vibrations of this kind may be called inextensional, or
flexural, and (§ 44) their frequencies are proportional to A, so that the
sounds become graver without limit as the thickness is reduced.

Vibrations of the one class may thus be considered to depend upon the term of
order h, and vibrations of the other class upon the term of order A , in the
expression for the potential energy.  In general both terms occur ; and it is
only in the limit that the separation into two classes becomes absolute. This
is a question which has sometimes presented difficulty. That in the case of
extensional vibrations the term in should be negligible in comparison with the
term in h seems reasonable enough. But is it permissible in dealing with the
other class of vibrations to omit the term in h while retaining the term in

The question may be illustrated by consideration of a statical problem. It is a
general mechanical principle (§ 74) that, if given 'displacements (not
sufficient by themselves to determine the configuration) be produced in a
system originally in equilibrium by forces of corresponding types, the
resulting deformation is determined by the condition that the potential energy
shall be as small as possible. Apply this principle to the case of an elastic
shell, the given displacements being such as not of themselves to involve a
stretching of the middle surface. The resulting deformation will, in general,
include both stretching and bending, and any expression for the energy will be
of the form

Ah (extension) + Bh (bending) (1).

This energy is to be ;is small as possible. Hence, when the thickness is
diminished without limit, the actual displacement will be one of pure bending,
if such there be, consistent with the given conditions.

At first sight it may well appear strange that of the two terms the one
proportional to the cube of the thickness is to be retained, while that
proportional to the first power may be neglected. The fact, however, is that
the large potential energy that would accompany any stretching of the middle
surface is the very reason why such stretching does not occur. The comparative
largeness of the coefficient (proportional to K) is more than neutralized by
the smallness of the stretching itself, to the square of which the energy is
proportional.

An example may be taken from the case of a rod, clamped at one end A, and
deflected by a lateral force; it is required to trace the effect of constantly
increasing stiffness of the part included between A and a neighbouring point B.
In the limit we may regard the rod as clamped at B, and neglect the energy of
the part JJ3, in spite of, or rather in consequence of, its infinite stiffness.

It would thus be a mistake to regard the omission of the term in h as
especially mysterious. In any case of a constraint which is supposed to be
gradually introduced (§ 92 cr), the vibrations tend to arrange themselves into
two classes, in one of which the constraint is observed, while in the other, in
which the constraint is violated, the frequencies increase without limit. The
analogy with the shell of gradually diminishing thickness is complete if we
suppose that at the same time the elastic constants are increased in such a
manner that the resistance to bending remains unchanged. The resistance to
extension then becomes infinite, and in the limit one class of vibrations is
purely inextensional or flexural.

In the investigation which we are about to give of the vibrations of a
cylindrical shell, the extensional and the inextensional classes will be
considered separately. It would apparently be more direct to establish in the
first instance a general expression for the potential energy complete as far as
the term in from which the whole theory might be deduced.  Such an expression
would involve the extensions and the curvatures of the middle surface. It
appears, however, that this method is difficult of application, inasmuch as the
potential energy (correct to A) does not depend only upon the above-mentioned
quantities, but also upon the manner of application of the normal forces, which
are in general implied in the existence of middle surface extensions.
\footnote{On the Uniform Deformation in Two Dimensional of a Cylindrical Shell, with
Application to the General Theory of Deformation of Thin Shells. Proc. Math,
Soc., vol. XX. p. 372, 1889.}

\section{235 c} % page 397 (411/495)
235 c. The first question to be considered is the expression of the conditions
that the middle surface remain unextended, or if these conditions be violated,
to find the values of the extensions in terms of the displacements of the
various points of the surface.

We will suppose in the first instance merely that the surface is of revolution,
and that a point is determined by cylindrical coordinates si, r, f. After
deformation the coordinates of the above point become z + Sz, r + Sr, f + Sl
respectively. If ds denote an element of arc traced upon the surface,

(ds -h (i8s)- = (dz -h dSzf -h (r + Sr)- (df  -h (iS )- + (dr -h dS ), so that

ds dSs = dz dhz -t- r~dfl dSj) + rSr (d )- + dr dSr (1).

In this we regard z and \(\phi\) as independent variables, so that, for example,
\(
d \delta z  =  d \delta z,
\)

while dr = (dr/dz) dr ,

in which by hypothesis \(dr/d \phi = 0\). Accordingly
\[
d \delta s (dzy (dhz dr dhr (ds) dz ds (dby f  dhf .
\]
(cZs)' ( d dz dij) (dhz dS dr dSr)

(ds)- d f dz  dz d(l .(2).  in which dSs ds represents the extension of the
element ds. If there be no extension of any arc traced upon the surface, (2)
must vanish independently of any relations bet een dz and d. Hence dhz dr dhr
dz .(3), O).  d2 , dhtb dr dhr From these, by elimination of hr.

dhz dr d dz V df) dhz dhd) dr d -h4 ~df  dz d(jr

and again, by elimination of hz.

,dhd dV d=80 dzV dz)' dz d(pr

% page 399
If the distribution of thickness and the form of the boundary or boundaries be
symmetrical with respect to the axis, the normal functions of the system are
to be found by assuming to be proportional to cos s0, or sin 50. The equation
for S0 may then be put into the form

It will be seen that the conditions of inextension go a long way towards
determining the form of the normal functions.

The simplest application is to the case of an cylinder for which r is constant,
equal say to a. Thus (3), (4), (5), (7) become simply dz. By (9), if S0 oc cos
50, we may take

a80 = (Aa- BgZ) cos sif (10),

and then, by (S), Sr = s (Aga 4- B z) sin 50 . (11),

= - 5 iia sin 50 (12).  Corresponding terms, with fresh arbitrary constants,
obtained by writing 50+i7r for s0, may of course bo added. If Bg = 0, the
displacement is in two dimensions only (§ 233).

If an inextensible disc be attached to the cylinder at z = 0, so as to form a
kind of cup, the displacements 8r and 80 must vanish for that value of z,
exception being made of the case 5=1. Hence Ag = 0, and

a 80 = 5,5 cos 50, 8r = 55,5 sin 50, 85 = - 5 5,a sin 50... (13).

Again, in the case of a cone, for which r = tan 7 . 5, the equations (3), (4),
(5), (7) become

dz . dSr  dS f

+ tan 7 =0, 5 tan 7 ~ 4- 8r = 0 dSz d80 dSr (15).  If we take, as usual, oc cos
s f, we get as the solution of (15) B = (Ag + cos s (16),

and corresponding thereto

Br = s tan y (A z + -B,) sin 5 (17),

Sz = tan 7 Bg - s ( A z + 5 ) sin (18).

If the cone be complete up to the vertex at = 0, = 0, so that

Ag cos S (19),

Br = sAg r sin y (20).

8z = - sAg tan y r sin stf .'..(21 ).  For the cone and the cylinder, the
second term in the general equation (7) vanishes. Wc shall obtain a more
extensive class of soluble cases by supposing that the surface is such that

r

= constant (22), an equation which is satisfied by surfaces of the second
degree in general. If Z 7

we shall find r  = -  (24) ;

dz a

and thus (7) takes the form


if Bj oc cos(s) , and a is defined by or in the present case The solution of (25)
is

r~ dz.  a . a + z (26).

,(27).  w

The corresponding values of Sr and Sz are to be obtained from (4) and (5).  If
the surface be complete through the vertex z = a, the term multiplied by B must
disappear. Thus, omitting the constant multiplier, we may take whence, by (4),
(5), Sr sb (a - r)   .

  - rrr-i sin sd

a (a + (29); (0).  .(31).  5 , b-(a - z) .

Bz = (sz + a)~-~~ smsl

If we measure s' from the vertex, z = a - z, and we may write cos (Sf) Sr = .9r
S2 = -S = 6--' (s + l)a (34).  For the parabola, a and b are infinite, while
b‘‘ a = 2a, and '=4oV. Thus we may take
\footnote{On the Infinitesimal Bending of Surfaces of Revolution. Proc. Math.
Soc., Tol. XIII. p. 4, 1881.}

S = r cos S(f), Sr = .9r + sin s(p, S.3' = - 2 ( 4- 1 ) a' r' sin . .(3r).

We will now take into consideration the
important case of the sphere, for which in (23) h = a. Denoting by 6 th‘ angle
between the radius vector and the axis, we have z = a cos 0, r = a sin 0, and
thus from (29), (30), (31) Sl = cos sl tan  0 (3G),

hr a == s sin 8 sin 0 tan  0 (37),

izja = (1 + 8 cos 0) sin tan  0 (88).  The other terms of the complete
solution, corresponding to (28), are to be obtained by changing the sign ol‘ s.

In the above equations the displacements are resolved parallel and
perpendicular to the axis \(\theta = 0\). It would usually be more convenient
to resolve along the normal and the meridian. If the components in these
directions be denoted by w and \( \alpha \delta \theta \), we have
\[
w = \delta r \sin (\theta) + \delta z  \cos (\theta),
\]
\[
\alpha \delta \theta = \delta r \cos (\theta) - \delta z  \sin (\theta);
\]
So that altogether

Bf = oos sf Af tan  0 + cot  0

= - sin s sin 04, tan ifl - col  0 (40), Wjd = sin sf) Afis 4- cos 0) tau  0
-f i, is - cos 0) cot  0 . . .(41).

To the above may be added terms derived by writing \( s \phi + \pi/2\) for
\( s \phi\), and changing the arbitrary constants.
\section{235 d}% page 401 (415/495)
235 d. We now proceed to apply the equations of § 235 c to the principal
extensions of a cylindrical surface, with a view to the formation of the
expression for the potential energy. The axial and circumferential extensions
will denoted respectively by
\(\epsilon_1\),
\(\epsilon_2\), and the shear by
\( \varpi \).
The first of these is given by (2) § 235 c, if we
suppose that \( d\phi =0\), \(dz/ds=1\).  Since in the case of a cylinder
\(dr/dz=0\), we find
\[\epsilon_1=\]
(1) In like manner
\[\epsilon_2=\]
(2).  The value of the shear may
be arrived at by considering the difference of extensions for the two
diagonals of an infinitesimal square whose sides are dz and acZA. It is
\[ \varpi = \].
(3).  The next part of the problem, viz. the expression of the
potential energy by means of \(\epsilon_1\), \(\epsilon_2\), \(\varpi\),
appertains to the general theory of elasticity, and can only be treated here in
a cursory manner. But it may be convenient to give the leading steps of the
investigation, referring for further explanations to the treatises
Thomson and Tait and of Love. In the notation of the former
(Natural Philosophy, § (694) the general equations in three dimensions are

M is Young’s modulus, sigma is Poisson’s ratio, n is the constant of
rigidity, and (m- n/3) that of compressibility.

ma = S, nb = T, nc=U (4).
Me   P   (Q+R)
Mf = Q - (R+P) sigma
Mg   R   (P+Q)
where sigma = (m-n)/(2m)

The energy w, corresponding to unit of volume is given by
\[
 2iv = (m + n) (e- + - + p-)
\]
\[
+ 2 (HI - It) ( ft +rfe + ef) + II (a- + b- + c-) (7 ).
\]
In the application to a lamina, supposed parallel to the plane ay, we are to
take Ji = 0, S = 0, T = 0, so that

,; = -0- -  , ( =0. 6=0 S).  Thus in terms of the extensions parallel to
, and if the shear c, ve get

f . . ai - n . . , , ,

W; = i J0--f (e+  + (it).  w=n

(9) This is the energy reckoned per unit of volume. In order to
adapt the expression to ur purposes, we must multiply it by the thickness
\(2h\)). Hence as the energy ier unit, area of a shell of thickness 26, we may
take in the notation adopted at the commencement of this section,
\[
 2nh \left( \epsilon_1^2 +\epsilon_2^2 + \frac{1}{2} \varpi^2 + \right)
\]
(10).  This expression may be applied to curved as well as to plane plates, for any
modification dm to curvature must involve higher powers ol It. The same is true
of the energy of bending.

\section{235 e}
235 e. We are now prepared for the investigation of the extensional vibrations
of an infinite cylindrical shell, assumed to be periodic with respect both to z
and \(\phi \).  It will be convenient to denote by single letters the
displacements parallel to \(z\), \(\phi \), \(r\); we take
\[\delta z=u , \delta \phi=v  , \delta r=w \]
(1), These functions are to be assumed proportional to the sines or cosines of
\(jz/a\), and \(s \phi \).  Various combinations may be made, of which an example
\footnote{Additions of \(\pi/2\) to \(s \phi \) or to \(jz/a\), or to both, may
of course be made at pleasure.} is
\[(u,v,w)=(U\cos(s \phi)\cos(jz/a), V \sin(s \phi )\sin(jz/a),
W\cos(s \phi )\sin(jz/a))\]
(2) so that (1),(2), (3) § 235 d,
\begin{align}
a \epsilon_1&=      -jU \cos(s \phi )  \sin(jz/a) \\
a \epsilon_2&= (W + sV) \cos(s \phi )  \cos(jz/a) \\
a \varpi    &= (-sU+jV) \sin(s \phi )  \cos(jz/a)
\end{align}
(3), (4), (5),
The potential energy per unit area is thus (10) § 235 d

 2'nJui~- j cos ir sin y3  + (li + sy + i sin sf cos -jz a ( -  fT + y F)=J (6).

Again, if \(\rho\) be the volume density, the kinetic energy per unit of area is

ph fdUy ,  , dVy sin 6 sin w j

(7).  In the integration of (6), (7) with respect to z and f,  is the
mean value of the square of each sine or cosine
\footnote{In the physical problem of a simple cylinder the range of integration
for 0 is from 0 to 2ir ; but mathematically we are not confined to one
revolution. We may conceive the shell to consist of several superposed
convolutions, and then  is not limited to be a whole number.}

We may then apply Lagrange’s method, regarding U, F, W as independent
generalized coordinates. If the type of vibration be cos joi, and p p n = i, the
resulting equations may be written

(2 (xV + 1) r + U-(2N -hi) F - 2Aj IF = 0. . .(S),

F+ 2(A + l)5Tr = 0...(9),
- 2Nj U- 2(X + l)sr+ (2 (X + 1 ) - k a l If = 0. . .(10), where JV' = ( 1 V

m -h n

The frequency equation is that expressing the evanescence of the determinant of
this triad of equations. On reduction it may be written
\footnote{Note on the Free Vibrations of an infinitely long Cylindrical Shell.
Proe Roy. Soc. vol. 45, p. 446, 1889.}


- ,s- (k a k') a - 2 (A -h 1) ( + + 1 )

-h 4 (2Ar 4- 1 ) j=J -f 4 (2A + 1 ) jV = 0 (1 2 V

These equations include of course the theory of the extensional vibrations of a
plane plate, for which a = x . In this application it is convergent to write
af = y, s a = 8, jja = 7. The displacements are then

u= U cos fiy cos 7 , v = V sin 0y sin yz, w = W cos fiy sin yz

(13)

When a is made infinite while 8, 7 remain
constant, the equations (10), (8), (0) ultimately assume the form = 0. and

(2 (.V+ 1) 7  + - k ) U (2N + 1) 7 SF = 0 ... (14),

- ( 2 JV + 1 ) 7 17 + ; 7 + 2 ( + 1 ) yd- A-- 1 1 - 5 ) :

and the determinantal equation (12) becomes

A-.  k‘  2 + 1) (r + yS- = 0 (16).

In (16), as was to be expected, k- appears as a function of (yS- + 7-).
The first root Ir = 0 relates to flexural vibrations, not here repeated
The second root is

Ar = yd -fr (17), or J- = ()9 +7 ) (18) At the same time (14) gives

7f ~y3r=0 (19).

These vibrations involve only a shearing of the plate in its own plane. For
example, if 7 = 0, the vibration may be represented by

n = cos ySy cos pt, v = 0, ; = 0 (20).

The third root of (16)

k =2(N+lHfi‘ + r) J'' jfi + Y) (21)

m -H n gives p- = (22).

The corresponding relation between U and V is

y3l7 + 7F=0 (23).

A simple example of this case is given by supposing in (13), (23). 8 = 0. We
may take

u = cos yz cos ptt V = 0, w = 0 (24),

the motion being in one dimension.  Reverting to the cylinder we will consider
in detail a few particular cases of importance. The first arises when j = 0,
that is, when the vibrations are independent of 2. The three equations (8),
(9), (10) then reduce to

(s- - A a ) 17 = 0 (25),

(2(JV'-hl)fi -fc a )F+ 2(7)7+ l)sW = 0 (26),

 2(i)7+l)F+(2(i)r + l)-A a ) 1F = 0 (27); and they may be satisfied in two ways.
First let F= = 0: then U may be finite, provided

= 0 (28).

The corresponding type for n is

u s cos cos pf ( 29),

where 2 = - (30 Y

to explain this motion the material is sheared without dilatation of area or volume,
every generating line of the cylinder moving along its own length. The
frequency depends upon the circumferential wave-length, and not upon the
curvature of the cylinder.

The second kind of vibrations are those for which that the motion is strictly
in two dimensions. The elimination of the ratio VjW from (26), (27) gives

I'-a - 2 (J + 1) (1 -h s)) = 0 (31 ), as the frequency equation. The first root
is fc- = 0, indicating infinitely slow motion. The modes in question are
flexural, for which, according to the present reckoning, the potential energy
is evanescent. The corresponding relation between V and M is by (26)

W =0 (32),
giving in (3), (4). (-)),
i = 0, e-j = 0, w = 0.

The other root of (31) is

= 2 (iVT + 1 ) (1+ s ) (33) or 4mn 1 H- s

Tn + n arp while the relation between V and W is .  (34); V sW  (35).

The type of the motion may be taken to be

u 3= 0, sin cos pt, w = cos cos pt (36).

It will be observed that when s is very large, the flexural vibrations (32)
tend to become exclusively radial, and the extensional vibrations (35) tend to
become exclusively tangential.

Another important class of i brat ions are those which are characterized by
symmetry round the axis, for which accordingly s = 0. The general frequency
equation (12) reduces in this case to

k a (r - 2 (.Y+ 1) (  + D + 4 (2.V + 1) = 0

... 37).

Corresponding to the first root we have r = 0, Tr = 0, as is readily proved on
reference to the original equations (8), (10

with s = 0. The vibrations are the purely torsional ones represented by

u = 0, v = sin ( jf o ) cos pt, iv = 0 ( 3S),



where P‘ = (39i

The frequency depends upon the wave-length parallel to the axis, and not on the
radius of the cylinder.

The remaining roots of (37) correspond to motions for which F=0, or which take
place in planes passing through the axis.  The general character of these
vibrations may be illustrated by the case where j is small, so that the
wave-length is a largo multiple of the radius of the cylinder. We find
approximately from the quadratic which gives the remaining roots khi- - 2iVj +

(40).

i +1

(41).

The vibrations of (40) are almost purely radial, that actually vanishes, we
fall back iijwn

If we suppose

kya'‘ = 2(N'+l)

(42).

and P'= , 711 -f- 71 uy

(43)

obtainable from (33), (34) on introduction of the condition j = 0.  The type of
vibration is now

M = 0, V = 0, W = CK)Apt

(44).

1 This equation was given by Love in a memoir On the Small Five
Vibrational Deformations of a thin Elastic Shell, Phil Trans.  vol. 179
fl888 p. it’s, and also by Chree, Cambridge Phil. Trans, vol. xiv, p. 1887.

On the other hand, the vibrations of (41) are ultimately purely axial. As the
type we may take

lit  71

u = cos jz a . cos pt, v = 0, - g- j sin jz a cos pt... (45), where j  3i - n
nj m pa Now, if q denote Young s modulus, we have, § 214, q = n (3m - n) m,
(46).  BO that . j47).

Thus u satisfies the equation

d u  q d

dt p dz

which is the usual formula (§ 150) for the longitudinal vibrations of a rod,
the fact that the section is here a thin annulus not influencing the result to
this order of approximation.

Another particular case worthy of notice arises when 1, so that (12) assumes
the form

k CL ( c a - j - 1) - 2 (JV + 1) (j + 2)

+ 4j (A; a - ) 1) = 0...(48).

As we have already seen, if j be zero, one of the values of vanishes. If j be
small, the corresponding value of is of the order Equation (48) gives in this
case jr+1

(49);

or in terms of p and q,

II

(50).

 The type of vibration is

u = 0  1

V  sin  sin jzja . cos pt

(51).

w = - cos 0 sin jzfa . cos pt )

and corresponds to the flexural vibrations of a rod (§ 163). In (51) V
satisfies the equation

in which ia represents the square of the radius of gyration of the section of
the cylindrical shell about a diameter.

This discussion of particular cases may suffice. It is scarcely necessary to
add, in conclusion that the most general deformation of the middle surface can
be expressed by means of a series of such as are periodic with respect to z and
f so that the problem considered is really the most general small motion of an
infinite cylindrical shell.

The extensional vibrations of a cylinder of finite length have been considered
by Love in his Theory of Elasticity’ (1S93), where will also be found a full
investigation of the general equations of extensional deformation.

\section{235 f Inextensional deformations}
235 f. When a shell is deformed in such a manner that no line traced upon the
middle surface changes in length, the term of order h disappears from the
expression for the potential energy, and unless we are content to regard this
function as zero, a further approximation is necessary. In proceeding to this
the first remark that occurs is that the quality of inextension attaches only
to the central lamina. Consider, for example, a portion of a cylindrical shell,
which is bent so that the original curvature is increased.  It is evident that
while the middle lamina remains unextended, those laminae which lie externally
must be stretched, and those that lie internally must be contracted. The amount
of these stretchings and contractions is proportional in the first place to the
distance from the middle surface, and in the second place to the change of
curvature which the middle surface undergoes. The potential energy of bending
is thus a question of the curvatures of the middle surface. Displacements of
translation or rotation, such as a rigid body is capable of, may be disregarded.

In order to take the question in its simplest form, let us refer, the original
surface to the normal and principal tangents at any point P as axes of
coordinates, and let us suppose that after deformation the lines in the sheet
originally coincident with the principal tangents are brought back
(if necessary) so as to occupy the same positions as at first. The possibility of
this will be apparent when it is remembered that, in virtue of the inextension
of the sheet, the angles of intersections of all lines traced

 Also Phil. Trans. vol. 179 a, 1888.  upon it remain unaltered. The equation of
the original surface in the neighbourhood of the point being that of the
deformed surface may be written 1 f V' (2).  In strictness (p + Spl)
(p.j + Sp) are the curvatures of the sections made by the planes a:, y ; but since the
principal curvatures are a maximum and a minimum they represent in general
with sufficient accuracy the new principal curvatures, although these are to be
found in slightly different planes. The condition of inextension shews that
points which have the same r, y in (1) and (2) are corresponding points ; and
by Gauss s theorem it is further necessary that

 '' + fe.O (3).

pi Pi It thus appears that the energy of bending will depend in general upon
two quantities, one giving the alterations of principal curvature, and the
other t depending upon the shift (in the material) of the principal planes.
The case of a spherical surface is in some respects exceptional.  Previously to
the bending there are no planes marked out as principal planes, and thus the
position of these planes after bending is of no consequence. The energy depends
only upon the alterations of principal curvature, and these by Gauss's theorem
are equal and opposite, so that, if a denote the radius of the sphere, the new
principal radii are a + a - 8p. If the equation of the deformed surface be 2r =
Aa + iBxy + Oif (4),

we have (a + ipY + (a - Sp) = .4 + G,

(a + Sp) (a -8p) = 4C- 5-;

so that 8 i (4 - C) + 5- (-5).  We have now to express the elongations of the
various laminas of a shell when bent, and we will begin with the case where t =
0, that is, when the principal planes of curvature remain unchanged.  It is
evident that in this case the shear c vanishes, and we have to deal only with
the elongations e and parallel to the axes, §235 d.  In the section by the
plane of r , let s, s' denote corresponding infinitely small arcs of the
middle surface and of a lamina distant h from it. If yr be the angle between
the terminal normals, s = piyjr, s' = (pi -f h) yjr, s -s = h yjr. lu the
bending, which loaves s unchanged,

hs =hSylr = hsh( pi).  Hence e = Ss'.s' = i S( 1

and in like manner f=hh(l p..). Thus for the energy U per unit aim we have 7U
4- n  p, Pa and on integration over the whole thickness of the shell (2 ) fs
Vs yj m

S iV p, V pj m4-n p, p.,, This conclusion may be applied at once, so as to
give the result applicable to a spherical shell ; for, since the original
principal planes are arbitrary, they can be taken so as to coincide with the
principal planes after bending. Thus t = 0 ; and by Gauss s tb core in

8(l pi)-h8(l pa) = 0,

so that

where 8 (1 p) denotes the change of principal curvature. Since e = - , 7 = 0,
the various lamina are simply sheared, and that in proportion to their
distance from the middle surface The energy is thus a function of the constant
of rigidity only.

The result (() is applicable directly to the plane j)late ; but this case is
peculiar in that, on account of the infinitude of pi , p, (3) is satisfied
without any relation between Spi and Sp.j. Thus for a plane plate

jj inh f 1 1 - n  1 1

U = - q i a d h '

3 (pi Pa' m n pi Pa

where l p , l p, are the two independent principal curvatures after

bending in A will be found to agree with the value (2) § 214, expressed in
a different notation.  We have thus far considered t to vanish; and it remains
to investigate the effect of the deformations expressed by 8z = Ta;y = iT(f-if)
(9), where f, i; relate to new axes inclined at 45' to those of x, y. The
curvatures defined by (9) are in the planes of f, ri, and are equal in
numerical value and opposite in sign. The elongations in these directions for
any lamina within the thickness of the shell are At,
- At, and the corresponding energy
(as in the case of the sphere just considered) takes the form
\[ U' = 4 n h^2 \tau^2/3\]
\todo[inline]{exponent of h is smudged}
(10)
This energy is to be added
\footnote{There are clearly no terms involving the products of \(\tau \)
with the changes of principal curvature
\(\delta \rho_1^{-1}\), \(\delta \rho_2^{-1}\); for a change
in the sign of \(\tau \) can have no influence upon the energy of the
deformation defined by (2).}
to that already found in (6) ; and we get finally as the complete expression of
the energy, when the deformation is such that the middle surface is unextended.
We may interpret t by means of the angle x through which the principal planes
are shifted; thus (12).

\section{235 g Application to Cylinder}
235 g. We will now proceed with the calculation of the
potential energy involved in the bending of a cylindrical shell.  The problem
before us is the expression of the changes of principal curvature and the
shifts of the principal planes at any point
\(P (z, \phi)\)
of the cylinder in terms
of the displacements
u, v, w. As in § 235 f, take as fixed coordinate axes the
principal tangents and normal to the undisturbed cylinder at the point P, the
axis of x being parallel to that of the cylinder, that of y tangential to the
circular section, and that of \(\zeta \)  normal, measured inwards. If, as it will be
convenient to do, we measure z and
\(\phi \) from the point P, we may express the
undisturbed coordinates of a material point Q In the neighbourhood of P, by
\[ x=z, y= a \phi ,  \zeta = \frac{1}{2} a \phi^2 \]
(1) During the displacement the coordinates of Q will receive the increments
n, w sin cos - m; cos + v sin  ; so that after displacement a: = j + z , y = +
wf -I- v (I f 0-) + vtp ;
or, if w, V, w be expanded in powers of the small quantities z, 0, du dll .
(2).
, , dn dv ,
= + 5
(3),
V 1 JO dw ,  ,
d w
dZtT -h d: w j dHv
w.  V if , , . being the values of i , v at the point P.

These equations give the coordinates of the various points of the deformed
sheet. We have now to suppose the sheet moved as a rigid body so as to restore
the position (as far as the first power of small quantities is concerned) of
points infinitely near P. A purely translatory motion by which the displaced P
is brought back to its original position will be expressed by the simple
omission in (2), (o), (4) of the terms I, iVo respectively, which are
independent of f).
The effect of an arbitrary rotation is represented by the
additions to \(x\), \(y\), \(\xi\) respectively of
\(y \omega_3-\xi \omega_2\), \(\xi \omega_1- x \omega_3\)
\(x \omega_2- y\omega_1\); where for the present purpose
\(\omega_1\), \(\omega_2\), \(\omega_3\)
are small quantities of the order of the deformation, the
square of which is to be neglected throughout. If we make these additions to
(2), etc., substituting for x, y, fin the terms containing 6 their approximate
values, we find so far as the first powers of z, f

da . du Now, since the sheet is assumed to be unextended, it must be possible
so to determine cwi,  to this order =

f=0. Hence

dll - dn do .

a.

dzQ dv .
- r,. 4- (1(0 = 0.
The conditions of inextension are thus (if we drop the suffices as no longer required)
J -0 dz do dll dv which agree with (S) § 235 c.

Returning to (2), etc., as indicated by the addition of the translatory and
rotatory terms, we get

r =  + terms of 2nd order in z, f,

y = 4 1 . 1 1 - d-w .

 = i a - + i  zl

I d'W dv dv

or since by  d-w ds- = 0, and dvjd = ~w,

 , ... 1 ... driv d-w dv ,

t J rf =  + rf..

The equation of the deformed surface after transference is thus ; ....(6).  (1
da 1 dw 1 L 1 1 ‘ ''1 fi

 odrcd i  (20 2a=  2a'-d )

Comparing with (2) § 235 we see that

t. 1 n pi 1 1  l

 p,ttr''d =) a(dsd4)‘

so that by (11) § 235

rr 4 i f 9a 1  dv dhv - jj 4 i ( VI If d y dv dHv

3a‘ (m + n a  d v Idr dzd (8).
This is the potential energy of bending reckoned per unit of area. It can, if
desired, be expressed by (o) entirely in terms of v
\footnote{From the general equations of Mr Love's memoir already cited a concordant
result may be obtained on introduction of the appropriate conditions.}

We will now apply (8) to calculate the whole potential energy of a complete
cylinder, bounded by plane edges r=± , and of thickness which, if variable at
all, is a function of s only. Since u, V, w are periodic when  increases by
 27r their most general expression in accordance with (5) is
[compare (1O), etc., § 235 c]
\begin{align}
v &= \sum (A +  B ) cos s'0 - ( A + B)  sin ( 0), \\
w &= \sum  s( A + B) sin s + s (A+B) cos  0 \\
u &= \sum - s (fi - BJa cos sf)
\end{align}
(11), in which the summation extends to all
integral values of s from 0 to \(\infty\). But the displacements corresponding to s
= 0, s = 1 are such as a rigid body might undergo, and involve no absorption
of energy. When the values of u, w are substituted in (8) all the terms
containing products of sines r cosines with different values of A vanish in
the integration with respect to 0, as do also those which contain cos sin ,s .
Accordingly to j. 47r  m Jo m n r sY (Aga + BfiZY + (Ag it -h B A  (A - 1 )'
4- V ) ,..(12).  Thus far we might consider h to be a function of r; but we
will now treat it as a constant. In the integration with respect to s the odd
powers of 3 will disappear, and wo get as the energy of the whole cylinder of
radius a, length 21, and thickness 2 , 17 jr j. 1   r a . I .  .(13), in which
A =2, 3, 4,....  The expression (13) for the potential energy suffices for the
solution of statical problems. As an example we will suppose that the cylinder
is compressed along a diameter by equal forces F, applied at the points z = 2u
=0, ff = 'rr, although it is true that so highly localised a force hardly comes
within the scope of

the investigation, in consequence of the stretchings of the middle surface,
which will occur in the immediate neighbourhood of the points of application
\footnote{
Whatever the curvature of the surface, an area upon
it may be taken so small as to behave like a plane, and therefore bend, in
violation of Gauss's condition, when subjected to a force which is so nearly
discontinuous that it varies sensibly within the area.
}


The work done upon the cylinder by the forces F during the hypothetical
displacement indicated by BAg, etc., will be by (10)

- (aSA  + ZiBBJ) (1 + cos sir),

so that the equations of equilibrium are dv = 0.  dA dv dA:

Thus for all values of s.  dv = - ( + cos Sir) saF,  - (1 -h cos sir) sz .
Ag = Bg= 0

and for odd values of a, Ag - Bg = 0.

But when s is even,

ms- 9sa F

m H- w  HmiliH (s- - 1)

(wi + n 3(1 j  SnnhH - 1)-

and the displacement vj at any point (z, is given by (14),

(15): w = 2 (AJa + B-fz) cos 2 + 4 (Ag'a + B z) cos 4 -h . . .(16), where A ,
B., -d ,... are determined by (14), (15).

A further discussion of this solution will be found in the memoir
\footnote{Proc. Roy. Soc. vol. 45, p. 105, 1888.} from which
the preceding results have been taken.

We will now proceed with the calculation for the frequencies of vibration of
the complete cylindrical shell of length 2Z. If the volume-density
\footnote{This can scarcely be confused with the notation for the curvature in
the preceding parts of the investigation.}
be \(\rho \), we
have as the expression of the kinetic energy by means of (9), (10), (11),
T = (1/2) 2hp integral integral (u‘ + v -h w ) adifdz

= 2 pi p h l a  sum (a (1 + s ) Ag)  + B  (17).

From the expressions for V and T in
(13), (17l the typos and frequencies of vibration can bo at once deduced. The
fact that the squares, and not the products, of A B are involved, shews that
these quantities are really the normal coordinates of the vibrating system. If
.4,, or A,', vary ns cos jj t, we have P 7tni A  - ,vV iu + It pa s- -I- I
This is the equation for the free frequencies of vibration in twt dimensions, § 233.
For a given material, the frequency is proportional directly to the
thickness and inversely to the square of the diameter of the cylinder .
\footnote{There is nothing in these laws special to the cylinder. In the case
of similar shells of any form, vibrating by pure bending, the frequency will be as the
thicknesses and inversely as corresponding areas. If the similarity extend also
to the thickness, then the frequency is inversely as the linear dimension, in
accordance with the general law of Cauchy.}

In like manner if or vary as cos f, we find 3n + i

 ‘  ..2T1 + a 3ff If the cylinder be at all long in proportion to its
diameter, the difference between p and becomes very small. Approximately in
this case , , - 3tt- i'ln + n 1

or, if we take m = 2, s = 2.

, , , 39n .(20);

\section{235 h Frequency Equations, Complete Sphere}
235 h. Complete Sphere. We now pass on to the consideration of spherical
shells. The general theory of the extensional vibrations of a complete shell
has been given by Lamb \footnote{Proc. Lend. Math. Soc. xiv. p. 60, 1882},
but as the subject is of small importance from an
acoustical point of view, we shall limit our investigation to the very simple
case of symmetrical radial vibrations.

If w be the normal displacement, the lengths of all lines upon the middle
surface aie altered in the ratio (a- -7v):a. In calculating the potential
energy we may take in (10) § 235 d,

= wja, tu = 0 ;



that the energy per unit area is 3m - 71 W' 4n i
or fur the whole sphere

F- 47ra- . 4fnh lu + u a-

 7il - n 7V' Hi -f n f .(1).  Also lor the kinetic energy, if p denote the
volume density,

T- . 47ra . 2A  .....(2).

Accordingly if tv=W cos pt, we have

dP VI + '

as the equation for the frequency (piir).

As regards the general theory Prof. Lamb thus summarizes his results.
\begin{quotation}
The fundamental modes of vibration fall into two classes. In the modes of the First
Class, the motion at every point of the shell is wholly tangential. In the nth
species of this class, the lines of motion are the contour lines of a surface,
harmonic \(S_n\) (Ch. XVII.), and the amplitude of vibration at any point is
proportional to the value of
\( dS/d \epsilon \), where \(d \epsilon \) is the angle subtended at the
centre by a linear element drawn on the surface of the shell at right angles
to the contour line passing through the point. The frequency
\( p 2 \pi \) is determined by the equation
\[ k^2 a^2 = (n-1)(n + 2),\]
(i), where a is the radius of the shell, and
\( k^2 = p^2 \rho/n\), if \(\rho \) denote the density, and
\(n\) the rigidity, of the substance.

In the vibrations of the Second Class, the
motion is partly radial and partly tangential. In the nth species of this class
the amplitude of the radial component is proportional to S, a surface harmonic
of order n. The tangential component is everywhere at right angles to the
contour lines of the harmonic Sn on the surface of the shell, and its amplitude
is proportional to
\( \Lambda  dS_n/d \epsilon \),
where \(\Lambda \), is a certain constant, and \( d \epsilon \),
has the same meaning as before.
\end{quotation}
Prof. Lamb finds k-Vr - 4fy 2n (i + 1) 7 (ii),
where k retains its former meaning, and 7 = (1 + a) (l - a),
denoting Poisson’s ratio. Corresponding to each value of n there are two values of
k-a , given by the equation

k'‘a'‘ - k-a- ( '- + t -f 4) 7 n-q- a - +4 (ti- + - 2 7 = 0. ..(iu).

Of the two roots of this equation, one is  and the other 47. It appears, then,
from (ii) that the corresponding fundamental modes are of quite different
characters. The mode corresponding to the lower root is always the more important.

 When n = l, the values of k'a- are 0 and 7. The zero root corresponds to a
motion of translation of the shell as a whole parallel to the axis of the
harmonic Sj. In the other node the radial motion is proportional to cos 6,
where 6 is the co-latitude measured from the pole of ; the tangential motion
is along the meridian, and its amplitude
(measured in the direction of 0 increasing) is proportional to i sin 0.

 When ii = 2, the values of ko corresponding lo various values of a are given
by the following table: r 0

1 I 1  Iu r  I rr M20  l lTrt Ms: 1190 ;io70 i 4-391 j 4G01 1 752 1

The most interesting variety is that of the zonal harmonic.  Making

S = i(3cos- -- 1),
we see that the polar diameter of the shell alternately elongates and
contracts, whilst the equator simultaneously contracts and expands
respectively. In the mode corresponding to the lower root, the tangential
motion is towards the poles when the polar diameter is lengthening, and vice
versa.  The reverse is the case in the other mode. We can hence understand the
great difference in frequency.

Prof. Lamb calculates that a thin glass globe 20 cm in diameter should, in its
gravest mode, make about 0.351 vibrations per second.

As regards inextensional modes, their form has already been determined,
(39) etc. § 235 c, at least for the case where the bounding curve and the thickness
are symmetrical with respect to an axis, and it will further appear in the
course of the present investigation. What remains to be effected is the
calculation of the potential energy of bending corresponding thereto, as
dependent upon the alterations of curvature of the middle surface. The process
is similar to that followed in § 2' og for the case of the cylinder, and
consists in finding the equation of the deformed surface when referred to
rectangular axes in and perpendicular to the original surface.

The two systems of coordinates to be connected are the usual polar coordinates
r, 6, and rectangular coordinates a:, y, measured from the point P, or (a, o),
on the undisturbed sphere. Of these x is measured along the tangent to the
meridian, y along the tangent to the circle of latitudes, and f along the
normal inwards.

Since the origin of 0 is arbitrary, we may take it so that 00 = 0. The relation
between the two systems is then = r Sin(tf - + sin cos o(l - cos 0)) (4),

y = r sin0 (5), f - r cos (0 - 6 ) - sin 00 sin 0 (1 - cos 0)-h u . . .

If we suppose r = , these equations give the rectangular coordinates of the
point (a, 0, 0) on the undisturbed sphere.  We have next to imagine this point
displaced so that its polar coordinates become a + 8r, 0-1-50, 0 -+- 80, and to
substitute these values in (4), (5), (G), retaining only the first power of Sr,
50, 50.  Thus

a; = (a -h Sr) (- sin (0 - 0o) + sin 0 cos 0o ( 1 - cos 0))

+ aS0 (- cos (0 - 0o) + cos 0 cos 0o(l - cos 0)1 -I- ft 50 sin 0 cos 00 sin 0
(7),

 = (a -H 5) sin 0 sin 0

-t- a50 cos 0 sin 0 + ftS0 sin 0 cos 0 (8).  f s=o - (o + 8r) (cos(d - d,) -
sin  sin  (1 - cos ))

+ ahO (sin (jS - 6 ) + sin cos  (1 - cos  ))

+ sin 00 sin d sin  (O').

These equations give the coordinates of any point Q of the sphere after
displacement ; but we shall only need to apply them in the case where Q is in
the neighbourhood of P, or (a, Bg, 0), and then the higher powers of (B - Bg)
and  may be neglected.

In pursuance of our plan we have now to imagine the displaced and deformed
sphere to be brought back as a rigid body so that the parts about P occupy as
nearly as possible their former positions. We are thus in the first place to
omit from (7), (8), (9) the terms (involving B) which are independent of - f.

Further we must add to each equation respectively the terms which represent an
arbitrary rotation, viz.

j I i - fj cos ,
determining ci,, cki , , in such a manner that, so far as the first powers of
(0 - o), 0, there shall be (coincidence between the original and displaced
positions of the point Q.

If we omit all terms of the second order in (0 - and we get from (7) etc.  ic =
-a(0 - 0o)- Sro (0 - 0J

- a Sda+  ( - + ... (10),

y = a sin  0 + Sr, sin 0 . + ahd cos 6 . f

+ a sin 0t S o+ (0 - 0) + i + a  f cos 0,(0 - 0) (11),

+ aB0 (0 - 0o) -h aS0o 0 i,f (12), Sr , etc. representing the values
appropriate t  P, whore (0 - 0, ) and (f vanish. The translation of the
deformed surface necessary to bring P back to its original position is
represented by the omission of the terms included in square brackets. The
arbitrary rotation is represented by the additions respectively of a sin
0,t.(f- cia, a(0~~ 0)3, - a (0 - ) oio - a sin 0.,f). w, :

and thus the destruction of the terms of the first order requires that

Sr a + dS0 d0 = O (IH),

- dS0 df) -h sin cos  + sin 0 g3 = 0 (14) ;

sin 0 dSffd0 + cos  S -f aj = 0 (1 o),

(Sr a) sin 0 cos 5 + sin 0 dSfjdf = 0 ( lb) ;

- dS (r tt) d5 + 80-6)2 = 0

- dS (rfa)l df) + sin° 0 8 - sin 0 o), = 0 ( 1 )

the suffixes being omitted.

These six equations determine w,cb, w,
giving as the three conditions of inextension

Sr a -h dB6;dd = 0 (19),

dB0, d(f + s 6 dB4d0 =0 (20),

Bv;it + cot 0B0 + dBlldf = 0 (21 ).

From (19), (20), (21), by elimination of S ,



01', since s 0 d d0 = d d log tan 0, rf   dBi

V  0)  d log tan 0~~ '

dB(f  d  f n-x

dl d log tan 0 sm0)

From (24), (25) we see that both Bif and S sin satisfy an equation of the
second order of the same form, viz.

d-ff Oft

fT log tan j )-  dd'

If the material system be symmetrical about the axis,  is a periodic function
of 0, and can be expanded by Fourier's theorem in a series of sines and
cosines of  and its multiples. Moreover each term of the series must satisfy
the equations independently.  Thus, if u varies as cos s0, ( 26) becomes d(log
tau S) - 5  = 0, whence a A' Um  0 + B' cot  0 (28),
where A' and B' are independent of 0. If we take

S(f = cos S(f -4, tan  0 + jB, cot  0(29),

we get for the corresponding value of B0 from (24)

S sin 5 = - sin sf A g tan k0 Bg cot  0(30) ;

and thence from (21) Sr o = sin sf -4 (s + cos 0) tan  0 + 5 , (s - cos 0) cot
0 .(31), as in (39), (40), (41) § 235 c.  The second solution (in B,) may be
derived from the first t iir d, in two waves which are both worthy of
notice. The manner of derivation from (27) shews that it is sufficient to
alter the sign of s.  tan becoming cot 1 , sin becoming - sin sj while cos si
remains unchanged.  The other method depends upon the consideration that the
general solution must be similarly i-iated to the two poles. It is thus
legitimate to alter the first solution by writing throughout (ir-O) ui plar-e
of 0, changing at the same time the sign of B0.

If we suppose ; = 1 , we get

sin dBj = CtS  .4, + 7, -  - .) rH.

80 = - i sin tA i - Bi - i A i -f 0

= sin (f) (i4i + i ,)sin tfj.

The displacement proportional to (Ai - Bi) is a rotation of the whole surface
as a rigid body round the axis 0 k7r, ( = () ; and that proportional to
(-4,j)
represents a translation parallel to the axis 0 = 7r,  = 7r. The
complementary translation and rotation with respect to the  ‘st' axes is obtained
bv substituting i) + Jtt for f).

The two other motions possible without bending correspond to a zero value of
\(s\), and are readily obtained from the original equations (19), (20), (21).
They are a rotation round the axis 0 = 0, represented by

80 = 0, 8 = constant, 8  0,

and a displacement parallel to the same n prescribed by

= 0. f ( .  ') = 0. - - cot e se,

a0 s f(

or S = 0, 80 = 7 sin 0, 8r = - 7 cos 0

If the sphere be complete, the displacements just considered , and
corresponding to j = 0, J, are the only ones possible. For higher values of ft
we see from (31) that 8r is infinite at one or other pole, unless A, and B both
vanish. In accordance with Jellet’s theorem
\footnote{On the Properties of Inextensible Surface,
Irish Acad. Trans. vol. 22, p 179, 1805.}
the complete sphere is incapable of bending.

If neither pole be included iu the actual surface, which for example we may
suppose bounded by circles of latitude, finite
values of both A and B are admissible, and therefore necessary for a complete
solution of the problem. But if, as would more often happen, one of the poles,
say 6 =0, is included, the constants B must be considered to vanish. Under
these circumstances the solution is

S(ff = Aa tan  0 cos

S0= - Aasm 0 0 Bin 8l I (32),

5 r = d , a (s + cos 0) tan J 0 sin sl )

to which is to be added that obtained by writing
\( s \phi + \pi/2\), and changing the arbitrary constant.

From (32) wo see that, along those meridians for which sinA‘ =0, the
displacement is tangential and in longitude only, while along the intermediate
meridians for which cos s(f) = 0, there is no displacement in longitude, but
one in latitude, and one normal to the surface of the sphere.

Along the equator d = 7r,

= Ag cos = - A sin sft Sr a = A si n s0,

so that the maximum displacements in latitude and longitude are equal.

Reverting now to the expressions for x, y, f in (7), (S), (9), with the
addition of the translatory and rotatory terms by which the deformed sphere is
brought back as nearly as possible to its original position, we kno r that so
far as the terms of the hrst order in ( - ft,) and  are concerned, they are
reduced to = - tt (ft - fto), y = a sin fto . f = 0 (33).

These approximations will suffice for the values of x and y ; but in the case
of f we require the expression complete so as to include all terms of the
second order. The calculation is straightforward. For any displacement such as
Sr in (9) we write y I I  1 dSr .  The additional rotatory terms are by (17),
(18)

L. IdBr) . f 1 dSr jl )

 - ; d .) + i ™ ft In these we are to retain only those terms in x, y,
which are of the second order and independent of S, so that we may write

x = af‘ sin cos e,u y = cos

In the complete expression for f as a quadratic function of (0 - 6o) and  thus
obtained, we substitute x and  from (3.S).  The final equation to the deformed
surface, after simplification by the aid of (19), (20), (21), may be written o
d-Sr)  ry j  I d-Br c(t 0 dSr

 2a sin a dffd(f  a df Br 1 d-Sr)

I Br cot 6 dBr df f 1   cot I 2a ( a sin

the suffixes being now unnecessary.

Taking the value of Br a from (32) we get

Sr Id-Sr - .

  - a 1 d Br a sin 0 d0dft + cos 5 d8r ai ui-0 d(f)  sin;= A f tan  0 cos Sip
    .( .30 ), Br cot 0 dBr a dO 1 d Br a sin 0 dtfr sin- 0   ‘ To obtain the
nmre complete solution corresponding to (31 ), we have only to add new terms,
multiplied by Bg, and derived from the above by changing the sign of s. As was
to be expected, the values in (35) and (37) are equal and opposite.

Introducing the values now found into (5) § 2 if we obtain as the square of
the change of principal curvature at any point

( p)  i sin tf

It should be remarked that, if either Ag or Bg vanish, (3S) is independent of
so that the change of principal curvature is the same for all points on a
circle of latitude, and that in any case (38) becomes independent of the
product ABg after integration round the circumference. The change of curvature
vanishes if s = 0, or  = 1, the displacement being that of which a rigid body
is capable.

Equations (35) etc. shew that along the meridians where Bp vanishes (cos(s) = 0)
the principal planes of curvature are the
meridian and its perpendicular, while along the meridians where Br vanishes,
the principal planes are inclined to the meridian at angles of 45.

The value of the square of the change of curvature obtained in (38) corresponds
to that assumed for the displacements in (29) etc., and for some purposes
needs to be generalised. We may add terms with coefficients A and B
corresponding to a change of s to (6  + i7r), and there is further to be
considered the summation with respect to s. Putting for brevity t in place of
tan d, we may take as the complete expression for S(l p)p, + 12  .  ((-4, -
Bgt- ) cos (sf) + (-4  - Bat'' ) cos (sl + tt) 1

When this is integrated with respect to ff round the entire circumference, all
products of the generalised coordinates Bg, Ag.Bg disappear,
so that (7)§ 235 f
we have as the expression for the potential energy of the surface included
between two parallels of latitude

V = 2w2 (. - 4-)11 sin 6 r

+ (39),

where

In the following applications to spherical surfaces where the pole  = 0 is
included, we may omit the terms in B and, if the thickness be constant, H may
be removed from under the integral sign. We have 2dt sin d = 2t

1+f- so that

Oi dB = i (1 + = 4 ~ .(41 ).

In the case of the hemisphere =1, and (41) assumes the value

(42).

 (=-)

Hence for a hemisphere uf uniform thickness

V= rHl' (s - ) (2 - 1) (A, + A.'-) (4.3).  If the extreme value of 6 be GO ,
instead of 90 , we get in place of (42)

4.  and V = Tr f S - s) (S,s- -f 4.9 - 3) (.4;- -i- A ). . .(4:)) These
expressions for V, in conjunction with (32), are sufficient for the solution
of statical problems, relative to the deformation o infinitely thin spherical
shells under the action of given impressed forces. Suppose, for example, that a
string of tension F connects the opposite points on the edge of a hemisphere,
represented by 0 = i7r, f = ir or tt, and that it is required to find the
deformation. It is evident from (32) that all the quantities -.4 vanish, and
that the work done by the impressed forces, corresponding t the deformation
Sil,, is
- BAgOS sin v T 4- sin F.

If s be odd this vanishes, and if s be even it is equal to
- iSAgOs sin hs7r,F.

Hence if s be odd Ag vanishes ; and by (43), if .9 be even, d VfdAg = ttH - 5)
(2ir - 1) = - 2as sin sw.F ; whence 2a F sin  sin 7ri (,9 --T(29=  l) (4G).
By (46) and (32) the deformation is completely determined.

If, to take a case in which the force is tangential, we suppose that the
hemisphere rests upon its pole with edge horizontal, and that a rod of weight W
is laid symmetrically along the diameter - tt, we find in like manner
a ir sin ,V7r irH (.9- - s) ( 2s - 1 ) (47) for all even values of js, and il = 0 for
all odd values of s.

We now proceed to evaluate the kinetic energy as defined by the formula) (T)‘

in which r denotes the surface density, supposed to be uniform.

If we take the complete value of Sf from (29), as supplemented by the terms in
A , BJ, we have

 (-4.'  + S,'t~‘).

When this expression is squared and integrated with respect to f round the
entire circumference, all products of letters with a different suffix, and all
products of dashed and undashed letters even with the same suffix, will
disappear. Hence replacing cos rsf etc. by the mean value i , we may take

sin e i sin dl(A + J; ) P

+ i sin' e s (4 + b;-) t- + sin e s (i,i. + A;b,').

The mean value (30) of
\(
(d \delta \theta/dt)^2
\)
is the same as that just written with the
substitution throughout of - B for B, so that we may take + sin  0 2 (4 + B; )
(49),

as the mean available for our present purpose. In (49) the products of the
symbols have disappeared, and if the expression for the kinetic energy were as
yet fully formed, the coordinates would be shewn to be normal.  But Ave have
still to include that part of the kinetic energy dependent upon dhrjdt As the
mean value, applicable for our purpose, we have from (31)

( i s (ir + i;) (H + cos er v

+ j 2 (B;- +  ) (s - cos ey r

, +2 ( A,B, + A,'B,') ( - cos- ff) (oO).

The expressions (49) and (oO) have now to be added. If we set for brevity

j tan d (- cos (d) - 2 sin (d) sin dd=f(s) (111),

or putting = 1 + cos 6,

 () = p K - ) -t- 2a (s -t- 1) - a) (ic (52).  we get

T = ro-a (S () (i,- + + :2 (- s) (Ji,- + B,'-)

+ 2Sj(s - cos tf)8m 0dO(A,i + ..-i ,) (.3)

It will be seen that, while V in (39) is expressible by the squares only of the
coordinates, a like assertion cannot in general be made of T. Hence Ag, etc.
arc 7iot in general the normal coordinates. Nor could this have been expected.
If, for example, we take the case where the spherical surface is bounded by
two circles of latitude equidistant from the equator, symmetry shews that the
normal coordinates are, not A and i, but (A -- B) and (A - B). In this case (-
s) = (s ).

A verification of (53) may readily be obtained in the particular case of J = l,
the surface under consideration being the entire sphere. Dropping the dashed
letters, we get

= i W ( 4 (i, -h B,y + g (ii - BJ (54).

In this case the displacements are of the purely translatory and rotatory
types already discussed, and the correctness of (54) may be confirmed.

Whatever may be the position of the circles of latitude by which the surface is
bounded, the true types and periods of vibration are determined by the
application of Lagrange s method to (39), (53).

When one pole, e,g. 0 = 0, is included within the surface, the coordinates B
vanish, and Ag become the normal coordinates.  If we omit the dashed
letters, the expression for T becomes simply

T=iiraa lf(8)Ag (55).

From (43), (55) the frequencies of free vibrations for a hemisphere are
immediately obtainable. The equation for Ag is

aa f(s) i, + if ( - if) (2 - 1) A, = 0 (56) ;

so that, if Ag vary as cos pgt,

  va is) 9  As) if we introduce the value of -ff from (40V and express a by
means of the volume density p.  .(58).  In like manner for the saucer of 120',
from (44),

A ) (8s H- 4s - 3)

ra f(s) .

The values of (s) can be calculated without difficulty in the various crises.
Thus, for the hemisphere,

 (2)= f ip (4 - 4r + )(1 + 6a; 1 = 20 log 2 -12J = 1-52961,

 (3) = 57i - 80 log 2 = 1-88156,

 (4) = 200 log 2 - 136J = 2-29609, etc. ;

so that p.,= j),= 14 726, x28'462.  a Ja  arsj(r  a‘Ja

In experiment, it is the intervals between the various
tones with which we are most concerned. We lind is Pa = 2-8102, pjpi = 5-4316
(59).

In the case of glass bells, such as are used with air-pumps, the interval
between the two gravest tones is usually somewhat smaller ; the representative
fi-action being nearer to 2-5 than 2-8.

For the saucer of 120°, the lower limit of the integral in (52) is and we get
on calculation

f(2) = -12864, (3) = 054884, giving Pi = X 7-9947, pj= x 20'911,

   a rr aVo'

p.  P-2 = 2-6157.  The pitch of the two gravest tones is thus decidedly
higher than for the hemisphere, and the interval between them is less.

With reference to the theory of tuning bells, it may be worth while to consider
the effect of a small change in the angle, for t he case of a nearly
hemispherical bell. In general X- = 4fr s - sY f sin 0 tan  0 d0 Jo aV f tan
0 a -h cos 0Y -h 2 sin 6 sin 0d0

J n ... (60).  If  = + and denote the value of p, for the exact hemisphere,
we get from previous results Thus

= PrO - MBe) = 'r(l -‘2t)Sdl, shewing that an increase in the angle depresses
the pitch. Am to the interval between the two gravest tones, we get showing
that it increases with 0. This agrees with the results given above for = GO .

The fact that the form of the normal functions is independent of the
distribution of density and thickness, provided that they vary only with
latitude, allows us to calculate a great variety of cases, the difficulties
being merely those of simple integration. If we suppose that only a narrow belt
in co-latitude
\(\theta \)
has sufficient thickness to contribute sensibly to the
potential and kinetic energies, we have simply, instead of (60),

\[ p^2 = 4 H s + cos f 2 s J
\]
whence
\[ yG + 4 cos 0 - cos' 0 ) p., V 1 1 1 + G
\]
\[
a T s + cos f 2 s J
\]
whence \(h  yG + 4 cos 0 - cos 0 \).
The ratio varies very slowly from 3, when
\(\theta = 0 \)
to
\(2.954\) when
\(\theta = \pi/2 \).

If 2 h denote the thickness at any co-latitude
\(\theta \),
\( H \propto h^2 \),
\( sigma  \propto h \),
I have calculated the ratio of frequencies of the two gravest tones
of a hemisphere on the suppositions (1) that
\( h \propto \cos(\theta) \),
and (2) that
\( h \propto 1 + \cos(\theta) \).
The
formula used is that marked (60) with
\( H \) and
\( sigma \),
under the integral signs. In the first case,
\( p_3 : p_2 = 1.7942\), differing
greatly from the value for a uniform thickness.  On the second
more moderate supposition as to the law of thickness,
\[ p_3 : p_2 = 2.4591,  \quad p_4 : p_2 = 4.4837. \]
It would appear that the smallness of the interval between the gravest tones
of common glass bells is due in great measure to the thickness diminishing with
increasing
\(\theta \).

\todo[inline]{ exponent of tan is smudged }
It is worthy of notice that the curvature of deformation
\(\delta( \rho^{-1} ) \),
which by (38) varies as
\( \sin^{-2} (\theta )  \tan (\theta )  \)
vanishes at the pole for \( s=3 \) and higher values, but is
finite for \( s=2 \).

The present chapter has been derived very largely from various published
memoirs by the author
\footnote{
Proc. Lfmd. Math. Soc., XIII. p. 4, 1881; xx. p. 372, 1889; Proc. Roy. Soc.,
vol.  45, p. 106, 1888; vol. 45. p. 448, 1888.
}.
The methods have not escaped criticism, some of which,
however, is obviated by the remark that the theory does not profess to be
strictly applicable to shells of finite thickness, but only to the 'limiting
case when the thickness is infinitely small. When the thickness increases, it
may become necessary to take into account certain local perturbations which
occur in the immediate neighbourhood of a boundary, and which are of such a
nature as to involve extensions of the middle surface. The reader who wishes to
pursue this rather difficult question may refer to memoirs by
Love \footnote{Phil. Trans., 179(a). p. 491. 1888; Proceeding Roy. Soc., vol.
49, p. 100, 1891; Theory of Electricity, ch. XXI.} ,
Lamb \footnote{ Proceedings. Lon. Math. Soc., vol. xxi. p. 119, 1890.} and
Basset \footnote{Phil. Trans. 181(a), p. 433, 1890; Am. Math: Journal., vol.
xvi, p. 264. 1894.}.  From the point of view of the present chapter the matter
is perhaps not of great importance. For it seems clear that any extension that
may occur must be limited to a region of infinitely small area, and affects
neither the types nor the frequencies of vibration. The question of what
precisely happens close to a free edge may require further elucidation, but
this can hardly be expected from a theory of thin shells. At points whose
distance from the edge is of the same order as the thickness, the
characteristic properties of thin shells are likely to disappear.

